% ---
% id: meat-core/thm:reduction
% kind: thm
% title: "Reduction of the rate to the convex program"
% summary: >-
%   For a declaration that adds no constraints on the state (an empty
%   constraint list), in either source mode, the value h_low of the finite
%   convex program bounds the single-round rate from below. Every
%   step of the reduction is an equality or lowers the bound, so the computable
%   program is conservative.
% prominence: main
% status: implemented
% uses: [def:hqkd-list, thm:protocol-channel, lem:cond-more, foundations/def:renyi, thm:branch-formula, def:declaration, def:admissible-declaration]
% mentions: [cor:reduction-constrained, protocol-compiler/rem:block-restriction, prop:a-neutral, thm:accumulated-entropy]
% assumes: []
% refs:
%   - "arqand2025meat, footnote 17"
% provenance: newly-derived
% implements: certiqs_key.meat_core.objective.objective:objective_value
% lean:
% notes: "Step 1 proves the restriction protocol-compiler/rem:block-restriction states; the compiler states it as a consumer obligation, and no claim declares it."
% ---

Step~1 of the proof is where the block restriction is discharged. The
compiler states, as a consumer obligation it does not discharge, that this
module's minimisation over the sector-block-diagonal feasible set is sound
(\xref{protocol-compiler/rem:block-restriction}); Step~1, the dephasing of an arbitrary
feasible state by the \emph{joint} pinching along the measured sectors and
along Alice's declared sectors, is that statement's proof. It is proved under
exactly the typing of Definition~\ref{def:declaration}\,(DX2), which types
Alice's measurement slot, $\tau_A$ and every $\Theta_k$ against a declared
factor. The restriction along $\mathbb{T}$ is licensed by this step, and
Proposition~\ref{prop:a-neutral} adds that it costs nothing. Corollary~\ref{cor:reduction-constrained} extends the theorem to
non-empty constraint lists.

\begin{theorem}\label{thm:reduction}
Let $\Decl$ satisfy Definition~\ref{def:admissible-declaration} with $\CList=\varnothing$,
in either source mode, and let $\alpha\in(1,\infty)$. Then, with
$\hqkd=\hqkd(\varnothing)$ the single-round rate of
Definition~\ref{def:hqkd-list} at the empty list,
\begin{equation}\label{eq:c8}
  \hqkd\;\ge\;\hlow .
\end{equation}
\end{theorem}

\begin{proof}
\emph{Idea.} Five steps take the bracket of \eqref{eq:hqkd-list} at
$\CList=\varnothing$, the rate that Theorem~\ref{thm:accumulated-entropy}
writes out as \eqref{eq:hqkd-def}, at a state of its inner infimum to the objective at a point of $\Feas$: a joint dephasing that
changes no output (Step~1), a purification (Step~2), the replacement of the
up-arrow entropy by the down-arrow one (Step~3), the branch formula of
Theorem~\ref{thm:branch-formula} (Step~4) and a rearrangement of the
logarithms (Step~5). Every step of the proof is an equality or an inequality
whose direction lowers the bound (conservative); no step uses $\CList$.

Fix $\mathbf q\in\acSet$ and a feasible $\omega\in\Seq{A\Qsys\Etil}$ with
$\omega\models\Decl$ (\eqref{eq:margcond}); abbreviate
$\Phi\coloneqq\Phimap[\omega_{A\Qsys}]$. We
lower-bound the bracket of \eqref{eq:hqkd-list} at $\CList=\varnothing$ by
$O_{\alpha}(\mathbf q,\tilde\omega)$ for some $\tilde\omega\in\Feas$; taking
infima then yields \eqref{eq:c8}. The divergence term of
\eqref{eq:hqkd-list} is already the KL divergence
$\frac{\alpha}{\alpha-1}\sfD(\mathbf q\|\Phi)$ of the import, so no
divergence relaxation occurs anywhere in this chain.

\emph{Step 1 (joint sector dephasing; exact).} Let
$\mathcal{P}_{\mathbb{T}}(X)\coloneqq\sum_{t\in\mathbb{T}}\Pi^{A}_tX\Pi^{A}_t$ be
the pinching along Alice's declared sectors
(Definition~\ref{def:declaration}\,(DX1); the identity map when
$|\mathbb{T}|=1$), $\PinS$ the pinching along $\SecSet$, for a product
\tref{sector map} the tensor product of the
\tref[measured factor]{measured factors}' pinchings, and
$\omega'\coloneqq(\mathcal{P}_{\mathbb{T}}\otimes\PinS\otimes\idmap_{\Etil})(\omega)$.
Then $\omega'$ is feasible: both pinchings are CPTP, so $\omega'$ is a
state, and in trusted-source mode
$\omega'_A=\mathcal{P}_{\mathbb{T}}(\omega_A)
=\mathcal{P}_{\mathbb{T}}(\tau_A)=\tau_A$ because $\tau_A$ is block diagonal
along $\mathbb{T}$ (Definition~\ref{def:admissible-declaration}\,(iii)). Its
$A\Qsys$-marginal is block diagonal along the joint sectors, and
$\nu^{\omega'}=\nu^{\omega}$: by \eqref{eq:povm-form} the output depends on
$\omega$ only through the operators
$\tr_{A\Qsys}[(M_x\otimes F_v\otimes\idon{\Etil})\,\omega]$, and
$\tr_{A\Qsys}[(M_x\otimes F_v\otimes\idon{\Etil})\,\omega']
=\tr_{A\Qsys}[(\mathcal{P}_{\mathbb{T}}(M_x)\otimes\PinS(F_v)\otimes
\idon{\Etil})\,\omega]
=\tr_{A\Qsys}[(M_x\otimes F_v\otimes\idon{\Etil})\,\omega]$, because both
pinchings are self-adjoint for the trace pairing and act on the traced
systems, every $F_v$ is block diagonal along $\SecSet$ and every $M_x$ along
$\mathbb{T}$ (Definition~\ref{def:admissible-declaration}\,(iii)). \emph{The step is an
equality}: $\omega'$ is another feasible point with literally the same
output state, Eve's register included; no dilation of a pinching is handed
to her here, and the purification is Step~2. Hence every entropy term and $\Phi$ are unchanged.
(If $\CList\neq\varnothing$, every entry still holds at $\omega'$; this is
Corollary~\ref{cor:reduction-constrained}.)

\emph{Step 2 (purification; conservative).} Let
$\ket{\hat\omega}\in\Hs{A\Qsys\Etil\Etil_2}$ purify $\omega'$, with
$\Etil_2$ of dimension $\dim\Hs{A\Qsys\Etil}$. Enlarging the purifying
register is within the scope of the import
(\cite[footnote~17]{arqand2025meat}).
For each $\hat c\in\Chkey$ with $q(\hat c)>0$: if $P_{\hat c}(\omega')=0$
then $\mathbf q\not\ll\Phi$ and the divergence term of the bracket is
already $+\infty$, so the claimed lower bound holds trivially; assume
$P_{\hat c}>0$. Since
$\tr_{\Etil_2}$ maps $\nu^{\hat\omega}_{|\hat c}$ to
$\nu^{\omega'}_{|\hat c}$ (conditioning on the classical $\Chat$ commutes
with the partial trace), Lemma~\ref{lem:cond-more} gives
\begin{equation*}
  \Hsandup(Z|I\Etil)_{\nu^{\omega'}_{|\hat c}}
  \;\ge\;\Hsandup(Z|I\Etil\Etil_2)_{\nu^{\hat\omega}_{|\hat c}}.
\end{equation*}

\emph{Step 3 (arrow replacement; conservative).} By the arrow order of
\xref{foundations/def:renyi},
$\Hsandup(Z|I\Etil\Etil_2)_{\nu^{\hat\omega}_{|\hat c}}
\ge\Hsanddown(Z|I\Etil\Etil_2)_{\nu^{\hat\omega}_{|\hat c}}$.

\emph{Step 4 (branch formula; exact).} $\hat\omega$ is pure with
$\hat\omega_{A\Qsys}=\omega'_{A\Qsys}\eqqcolon\tilde\omega$
block diagonal along the joint sectors, and $\tilde\omega\in\Feas$: normalization is
preserved by every step ($\tr\tilde\omega=\tr\omega=1$), and in
trusted-source mode additionally
$\tilde\omega_A=\omega_A=\tau_A$, none of Steps 1--3 having touched the
$A\Qsys$-marginal except through the joint pinching of Step~1, which
fixes $\tau_A$. Theorem~\ref{thm:branch-formula}, applied with
purifying register $\Etil\Etil_2$, gives
\begin{equation*}
  \Hsanddown(Z|I\Etil\Etil_2)_{\nu^{\hat\omega}_{|\hat c}}
  =\frac{1}{1-\alpha}\log_2
  \frac{K_{\hat c}(\tilde\omega)}{P_{\hat c}(\tilde\omega)}
  =\frac{1}{1-\alpha}\log_2\frac{K_{\hat c}(\tilde\omega)}{P_{\hat c}},
\end{equation*}
using $P_{\hat c}(\tilde\omega)=P_{\hat c}(\omega')=P_{\hat c}$.

\emph{Step 5 (rearrangement; exact).} Collecting Steps 1--4, the bracket of
\eqref{eq:hqkd-list} is bounded from below by
\begin{equation*}
  \frac{\alpha}{\alpha-1}\sfD(\mathbf q\|\Phi)
  +\sum_{\hat c\in\Chkey}
  \frac{q(\hat c)}{1-\alpha}\log_2
  \frac{K_{\hat c}(\tilde\omega)}{P_{\hat c}} .
\end{equation*}
For each $\hat c\in\Chkey$ with $q(\hat c)>0$ and
$P_{\hat c},K_{\hat c}>0$,
\begin{equation}\label{eq:rearrange}
\begin{split}
  \frac{\alpha}{\alpha-1}\,q(\hat c)\log_2\frac{q(\hat c)}{P_{\hat c}}
  &+\frac{q(\hat c)}{1-\alpha}\log_2\frac{K_{\hat c}}{P_{\hat c}}\\
  &=q(\hat c)\log_2\frac{q(\hat c)}{P_{\hat c}}
  +\frac{q(\hat c)}{\alpha-1}
  \left[\log_2\frac{q(\hat c)}{P_{\hat c}}
  -\log_2\frac{K_{\hat c}}{P_{\hat c}}\right]\\
  &=q(\hat c)\log_2\frac{q(\hat c)}{P_{\hat c}}
  +\frac{q(\hat c)}{\alpha-1}\log_2\frac{q(\hat c)}{K_{\hat c}},
\end{split}
\end{equation}
an identity of real numbers. At the degenerate values the identity persists
in the extended sense: $q(\hat c)=0$ makes all terms $0$; $P_{\hat c}=0$
with $q(\hat c)>0$ makes both sides $+\infty$; $K_{\hat c}=0$ with
$q(\hat c)>0$ makes both sides $+\infty$ (using $K_{\hat c}\le P_{\hat c}$
from Theorem~\ref{thm:branch-formula}'s \emph{cone clause}, which is stated without the
hypothesis $P_{\hat c}>0$ that the degenerate cases here violate, so the
$+\infty$ appears on both sides through
the $K$ term on the right and through the combination on the left, where
$\frac{1}{1-\alpha}\log_2\frac{K}{P}\to+\infty$). Summing
\eqref{eq:rearrange} over $\hat c\in\Chkey$ and adding the
$\hat c\notin\Chkey$ part of $\sfD$ yields exactly
$O_{\alpha}(\mathbf q,\tilde\omega)\ge\hlow$ after taking the infimum over
$\Feas$ and then over $\acSet$.
\end{proof}
